% =====================================================================
%  PRÉAMBULE COMMUN — Carnet d'entraînement, Fiche 2 (Nomenclature)
%  (inséré physiquement dans chaque document pour qu'il soit autonome)
% =====================================================================
\documentclass[11pt,a4paper]{article}

% ---------- Encodage & langue ----------
\usepackage[utf8]{inputenc}
\usepackage[T1]{fontenc}
\usepackage{lmodern}
\usepackage[french]{babel}

% ---------- Mathématiques & chimie ----------
\usepackage{amsmath,amssymb,mathtools}
\usepackage{icomma}
\usepackage[version=4]{mhchem}
\usepackage{siunitx}
\sisetup{output-decimal-marker={,},per-mode=symbol}

% ---------- Graphismes & mise en page ----------
\usepackage{xcolor}
\usepackage{colortbl}
\usepackage{tikz}
\usepackage[most]{tcolorbox}
\usepackage{enumitem}
\usepackage{array}
\usepackage{booktabs}
\usepackage{multicol}
\usepackage{fancyhdr}
\usepackage{titlesec}
\usepackage[hidelinks]{hyperref}
\usetikzlibrary{arrows.meta,positioning,calc,shapes.geometric,shapes.misc,
  fit,backgrounds,decorations.pathreplacing}
\usepackage[a4paper,margin=2.2cm,headheight=28pt,headsep=10pt]{geometry}

% =====================================================================
%  PALETTE (charte « Chimie » — mauve/prune du cours)
% =====================================================================
\definecolor{primary}{RGB}{146,95,161}
\definecolor{primarydark}{RGB}{92,55,108}
\definecolor{defcol}{RGB}{0,114,178}
\definecolor{thmcol}{RGB}{0,140,120}
\definecolor{formulacol}{RGB}{198,93,9}
\definecolor{remarkcol}{RGB}{120,94,200}
\definecolor{attncol}{RGB}{200,40,55}
\definecolor{excol}{RGB}{0,135,90}
\definecolor{methcol}{RGB}{90,90,100}
\definecolor{metalcol}{RGB}{198,150,40}
\definecolor{nonmetcol}{RGB}{0,140,150}
\definecolor{oxycol}{RGB}{210,55,45}
\definecolor{hydcol}{RGB}{60,120,210}
\definecolor{catcol}{RGB}{200,70,120}
\definecolor{ancol}{RGB}{40,110,190}
\definecolor{saltcol}{RGB}{120,90,160}

% =====================================================================
%  STYLE COMMUN DES BOÎTES
% =====================================================================
\tcbset{
  boxbase/.style={enhanced,breakable,boxrule=0.9pt,arc=2.5pt,
     left=3mm,right=3mm,top=2.3mm,bottom=2.3mm,
     fonttitle=\bfseries,coltitle=white,
     toptitle=0.6mm,bottomtitle=0.6mm},
}

\newtcolorbox{definition}[1][]{%
  boxbase,colback=defcol!5,colframe=defcol,
  title={\textsc{Définition}\ifx\relax#1\relax\relax\else\textmd{ --- \itshape #1}\fi}}
\newtcolorbox{regle}[1][]{%
  boxbase,colback=thmcol!6,colframe=thmcol,
  title={\textsc{Règle}\ifx\relax#1\relax\relax\else\textmd{ --- \itshape #1}\fi}}
\newtcolorbox{formule}[1][]{%
  boxbase,colback=formulacol!6,colframe=formulacol,
  title={\textsc{Méthode-clé}\ifx\relax#1\relax\relax\else\textmd{ --- \itshape #1}\fi}}
\newtcolorbox{remarque}[1][]{%
  boxbase,colback=remarkcol!6,colframe=remarkcol,
  title={\textsc{Remarque}\ifx\relax#1\relax\relax\else\textmd{ : \itshape #1}\fi}}
\newtcolorbox{attention}[1][]{%
  boxbase,colback=attncol!6,colframe=attncol,
  title={\textsc{Attention}\ifx\relax#1\relax\relax\else\textmd{ : \itshape #1}\fi}}
\newtcolorbox{exemple}[1][]{%
  boxbase,colback=excol!5,colframe=excol,
  title={\textsc{Exemple}\ifx\relax#1\relax\relax\else\textmd{ : \itshape #1}\fi}}
\newtcolorbox{methode}[1][]{%
  boxbase,colback=methcol!7,colframe=methcol,sharp corners,
  title={\textsc{Méthode}\ifx\relax#1\relax\relax\else\textmd{ : \itshape #1}\fi}}
\newtcolorbox{astuce}[1][]{%
  boxbase,colback=primary!6,colframe=primary,
  title={\textsc{Astuce concours}\ifx\relax#1\relax\relax\else\textmd{ : \itshape #1}\fi}}

% ---- Exercice (numéroté) ----
\newtcolorbox[auto counter]{exercice}[1][]{%
  boxbase,colback=primary!4,colframe=primarydark,
  title={\textsc{Exercice}~\thetcbcounter\ifx\relax#1\relax\relax\else\textmd{ --- \itshape #1}\fi}}

% ---- Corrigé (numéroté) ----
\newtcolorbox[auto counter]{corrige}[1][]{%
  boxbase,colback=excol!4,colframe=excol,
  title={\textsc{Corrigé}~\thetcbcounter\ifx\relax#1\relax\relax\else\textmd{ --- \itshape #1}\fi}}

% =====================================================================
%  EN-TÊTES ET PIEDS DE PAGE
% =====================================================================
\pagestyle{fancy}
\fancyhf{}
\fancyhead[L]{\small\textcolor{primarydark}{\textbf{Fiche 2} — Nomenclature}}
\fancyhead[R]{\small\textcolor{primarydark}{\textsc{Carnet d'entraînement}}}
\fancyfoot[C]{\small\thepage}
\renewcommand{\headrulewidth}{0.8pt}
\renewcommand{\headrule}{\hbox to\headwidth{\color{primary}\leaders\hrule height \headrulewidth\hfill}}

\titleformat{\section}{\Large\bfseries\color{primarydark}}{\thesection}{0.6em}{}
\titleformat{\subsection}{\large\bfseries\color{primary}}{\thesubsection}{0.6em}{}
\titleformat{\subsubsection}{\normalsize\bfseries\color{primary!80!black}}{}{0em}{}

% ---- Bandeau de titre réutilisable : \bandeau{Thème}{Sous-titre} ----
\newcommand{\bandeau}[2]{%
  \begin{center}
  \begin{tikzpicture}
  \node[rectangle,rounded corners=7pt,fill=primary,text=white,
        minimum width=16.0cm,minimum height=1.7cm,align=center,inner sep=8pt]
    {{\Large\bfseries #1}\\[3pt]{\normalsize #2}};
  \end{tikzpicture}
  \end{center}
  \vspace{0.3cm}}
\begin{document}
\bandeau{Thème 1 — Moyen}{Ions, valences et nombre d'oxydation --- Corrigé}

\begin{corrige}[n.o.\ dans un oxyde neutre]
On pose $x$ = n.o.\ cherché, on écrit « $x\cdot(\text{nb atomes}) + (-\mathrm{II})\cdot(\text{nb }\ce{O}) = 0$ » :
\begin{itemize}[leftmargin=1.5em,topsep=1pt,itemsep=1pt]
  \item \ce{SO3} : $x + 3(-2)=0 \Rightarrow x=+\mathrm{VI}$. \quad (soufre $+\mathrm{VI}$)
  \item \ce{N2O5} : $2x + 5(-2)=0 \Rightarrow 2x=10 \Rightarrow x=+\mathrm{V}$. \quad (azote $+\mathrm{V}$)
  \item \ce{Cl2O5} : $2x + 5(-2)=0 \Rightarrow 2x=10 \Rightarrow x=+\mathrm{V}$. \quad (chlore $+\mathrm{V}$)
  \item \ce{CO2} : $x + 2(-2)=0 \Rightarrow x=+\mathrm{IV}$. \quad (carbone $+\mathrm{IV}$)
\end{itemize}
\end{corrige}

\begin{corrige}[n.o.\ dans un composé avec hydrogène]
\begin{itemize}[leftmargin=1.5em,topsep=1pt,itemsep=1pt]
  \item \ce{NH3} : $x + 3(+1)=0 \Rightarrow x=-\mathrm{III}$. \quad (azote $-\mathrm{III}$)
  \item \ce{H2S} : $2(+1) + x =0 \Rightarrow x=-\mathrm{II}$. \quad (soufre $-\mathrm{II}$)
  \item \ce{CH4} : $x + 4(+1)=0 \Rightarrow x=-\mathrm{IV}$. \quad (carbone $-\mathrm{IV}$)
  \item \ce{H2SO4} : $2(+1) + x + 4(-2)=0 \Rightarrow x = 8-2 = +\mathrm{VI}$. \quad (soufre $+\mathrm{VI}$)
\end{itemize}
Note : le soufre vaut $-\mathrm{II}$ dans \ce{H2S} mais $+\mathrm{VI}$ dans \ce{H2SO4} : le n.o.\ dépend
\emph{du composé}, pas seulement de l'élément.
\end{corrige}

\begin{corrige}[n.o.\ dans un ion polyatomique]
La somme des n.o.\ égale la charge de l'ion :
\begin{itemize}[leftmargin=1.5em,topsep=1pt,itemsep=1pt]
  \item \ce{MnO4-} : $x + 4(-2) = -1 \Rightarrow x = -1+8 = +\mathrm{VII}$. \quad (manganèse $+\mathrm{VII}$)
  \item \ce{Cr2O7^2-} : $2x + 7(-2) = -2 \Rightarrow 2x = -2+14 = 12 \Rightarrow x = +\mathrm{VI}$. \quad (chrome $+\mathrm{VI}$)
  \item \ce{SO4^2-} : $x + 4(-2) = -2 \Rightarrow x = -2+8 = +\mathrm{VI}$. \quad (soufre $+\mathrm{VI}$)
  \item \ce{NO3-} : $x + 3(-2) = -1 \Rightarrow x = -1+6 = +\mathrm{V}$. \quad (azote $+\mathrm{V}$)
\end{itemize}
\end{corrige}

\begin{corrige}[du n.o.\ au nombre d'électrons]
Électrons restants $= Z - n$ :
\begin{center}
\begin{tabular}{lcccc}
\toprule
Cation & \ce{Fe^3+} & \ce{Cu+} & \ce{Cu^2+} & \ce{Sn^4+} \\
\midrule
Électrons & $26-3=23$ & $29-1=28$ & $29-2=27$ & $50-4=46$ \\
\bottomrule
\end{tabular}
\end{center}
\end{corrige}

\begin{corrige}[un métal, deux degrés]
\begin{itemize}[leftmargin=1.5em,topsep=1pt,itemsep=1pt]
  \item \ce{FeO} : $x + (-2) = 0 \Rightarrow x = +\mathrm{II}$ $\Rightarrow$ \textbf{fer(II)}.
  \item \ce{Fe2O3} : $2x + 3(-2) = 0 \Rightarrow 2x = 6 \Rightarrow x = +\mathrm{III}$ $\Rightarrow$ \textbf{fer(III)}.
\end{itemize}
Le fer peut prendre \emph{plusieurs} valences (II ou III) : la même paire d'éléments donne deux composés
différents. Il faut donc préciser « fer(II) » ou « fer(III) » pour lever l'ambiguïté. Le sodium, lui, n'a
qu'\emph{une seule} valence ($+\mathrm{I}$) : \ce{Na2O} ne peut être que « oxyde de sodium », inutile
d'indiquer la valence.
\end{corrige}
\end{document}
